\scanpage{95}
% archive foliation: 94
\unsure{donne|trouve} l'isom
\[
\kappa_\alpha:T_0^{\otimes2}\simeq T_{\Sigma,t}
\]
\unsure{à} une courbe elliptique $E_\alpha$, \inserted{corr. à $\underline J,t$} munie
d'une \deleted{\unsure{rigidification}}\inserted{épinglage} de Legendre subordonné
à $\alpha$ (puisque $T_{E_\alpha,0}\simeq T_0\ \ldots$). Ceci
posé, pour une courbe elliptique quelconque
$E$ sur $S$, corr. à $\underline J,t$, les épinglages
de Legendre sur $E$ subordonnés à $\alpha$
forment (on l'a vu) les sections d'un
$\unsure{\boldsymbol{\mu}|\mu}_2$-torseur $\mathcal L_\alpha(E)$ sur $S$, et on voit
de suite qu'on a un isom de
$\unsure{\boldsymbol{\mu}|\mu}_2$-torseurs
% margin position: left, diagonally
\margin{c'est la courbe elliptique \uline{universelle}, correspondant
à $\underline J,t$ et munie d'un épinglage de Legendre
subordonné à $\alpha$\ldots}
\[
\deleted{\unsure{\mathcal L_\alpha(E)\longleftarrow}}
\]
\[
(57)\qquad
\underbrace{\underline{\operatorname{Isom}}_{S,J}(E_\alpha,E)}_{
\substack{\text{l'indice }J\text{ signifie}\\
\text{qu'on prend les iso.}\\
\text{compatibles avec}\\
\underline J\simeq{}_2E\simeq{}_2E_\alpha}}
\xrightarrow{\sim}\mathcal L_\alpha(E)
\]
\[
u\longmapsto u(\kappa_\alpha)
\]
Si, au lieu d'une section\inserted{$\alpha$} de $\unsure{\boldsymbol{\mu}|\mu}_4^*\wedge\underline\omega_{\underline J}$, on
disposait d'une section $\omega$ de $\underline\omega_{\underline J}$ lui-même,
on pourrait encore définir une courbe
elliptique $E_\omega$ canonique corr. à $\underline J,t$, et
munie d'un ``ép. de Legendre ``naïf'' ''
(i.e. d'un iso. $T_{E,0}\simeq\underline O_\Sigma(1)(t)$ tel que\ldots), et
